%%% The macro file where all style boxes are defined. 
%%% We should make sure we get this right before we distribute more widely.
%%% We should also clean this up a little. It's a working mess that I threw together as I needed
%%% Compiled by Thomas Wong, 2020


\usepackage{amsmath,amssymb,amsthm}%,import,dashrule}
\usepackage{graphicx}

\usepackage{multicol}
\usepackage{pifont}
\usepackage{amsfonts}
\usepackage{aliascnt}
\usepackage{fancyhdr}
\setlength{\headheight}{15pt}
\usepackage{enumitem}
\usepackage{lmodern,textcomp}%fonts

\usepackage{wasysym}


%\usepackage{gensymb} %for degree symbol

\usepackage{tikz}
\usetikzlibrary{fit,shapes}
\usetikzlibrary{plotmarks}
\usetikzlibrary{matrix,decorations.pathreplacing}
\usetikzlibrary{automata,positioning, calc}
\usepackage{pgfplots}

\usepackage[RPvoltages]{circuitikz} %drawing circuits
\usepackage[hypcap=false]{caption} %captioning figures
\renewcommand{\thefigure}{\thesection.\arabic{figure}}

\usepackage[titles]{tocloft}
\usepackage[paper=letterpaper,left=25mm,right=25mm,top=3cm,bottom=25mm]{geometry}
\usepackage{parskip}

\usepackage{colortbl}
\usepackage{color}   %May be necessary if you want to color links

\newcommand{\diag}{\operatorname{diag}}

\newcommand{\blue}{blue}
\newcommand{\red}{red}
\newcommand{\brown}{brown}
\newcommand{\green}{green}
\newcommand{\yellow}{yellow}
\newcommand{\gray}{gray!30}
\newcommand{\darkgray}{gray!60}
\newcommand{\black}{black}
\newcommand{\graphcolour}{red}
\newcommand{\white}{white}



\colorlet{lblue}{\blue!45!white}
\colorlet{lred}{\red!50!white}
\colorlet{lblue1}{\blue!10!white}
\colorlet{lred1}{\red!10!white}
\colorlet{lblue2}{\blue!45!white!50!\black}

% \newcommand{\vv}[1]{\underline{#1}}
 \newcommand{\vi}{\underline{i}}
 \newcommand{\vj}{\underline{j}}
 \newcommand{\vk}{\underline{k}}

%\newcommand{\vv}[1]{{\mathbf{#1}}}
%\newcommand{\vi}{\vv{i}}
%\newcommand{\vj}{\vv{j}}
%\newcommand{\vk}{\vv{k}}

\usepackage{hyperref}
\hypersetup{
    colorlinks=true, %set true if you want colored links
    linktoc=all,     %set to all if you want both sections and subsections linked
    linkcolor=lblue!50!black,  %choose some color if you want links to stand out
		filecolor=lblue!50!black,
		urlcolor=lblue!50!black,
		pdfnewwindow=true
}

\usepackage{polynom}
\usepackage[]{geometry}
\usepackage{array}
\usepackage{pdflscape}
\usepackage{rotating}

\usepackage{tcolorbox}
\usepackage{varwidth}
\tcbuselibrary{skins,breakable}
\usetikzlibrary{shadings,shadows}

\usepackage{chngcntr}%changes counter
\counterwithin*{figure}{section}

%creating QR codes
\usepackage{marginnote}
\usepackage{qrcode}

%writing code
\usepackage{verbatim}
%\usepackage{minted} % for code snippets
%\usepackage{listings}


% \lstset{
%   basicstyle=\ttfamily\small,
%   keywordstyle=\color{blue},
%   commentstyle=\color{gray},
%   stringstyle=\color{orange},
%   breaklines=true,
%   frame=single,
%   showstringspaces=false
% }

\renewcommand{\familydefault}{\sfdefault}

%\setcounter{section}{-1}
%\setlength{\cftbeforechapskip}{2ex}
\setlength{\cftbeforesecskip}{0.1ex}
\setcounter{tocdepth}{2}
\pagestyle{fancyplain}
\numberwithin{equation}{section}
 
\theoremstyle{definition}
\newtheorem{cou}{counting}[section]

\newaliascnt{df}{cou}
\newaliascnt{thm}{cou}
\newaliascnt{rem}{cou}
\newaliascnt{example}{cou}
\newaliascnt{exercise}{cou}

%\newtheorem{df}[df]{Definition}
%\newtheorem{thm}[thm]{Theorem}
%\newtheorem{example}[example]{Example}
%\newtheorem{exercise}[exercise]{Exercise}

%\newtcolorbox[use counter=thm,number within=section]{thm}[2][]{%
%fonttitle=\bfseries,
%title= Theorem \thethm\; #2,#1}

\newtcolorbox[use counter=exercise,number within=section]{exercise}[2][]{%
fonttitle=\bfseries,
colbacktitle=lblue!60!black,
title= Exercise \theexercise\; #2,#1}

\newtcolorbox[use counter=rem,number within=section]{rem}[2][]{enhanced,
colbacktitle=\blue,
boxrule=0.2mm,
attach boxed title to top left={xshift=1cm,yshift*=1mm-\tcboxedtitleheight},
varwidth boxed title*=-3cm,
boxed title style={frame code={
\path[fill=tcbcolback!30!black]
([yshift=-1mm,xshift=-1mm]frame.north west)
arc[start angle=0,end angle=180,radius=1mm]
([yshift=-1mm,xshift=1mm]frame.north east)
arc[start angle=180,end angle=0,radius=1mm];
\path[left color=tcbcolback!60!black,right color=tcbcolback!60!black,middle color=tcbcolback!60!black]
([xshift=-2mm]frame.north west) -- ([xshift=2mm]frame.north east)
[rounded corners=1mm]-- ([xshift=1mm,yshift=-1mm]frame.north east)
-- (frame.south east) -- (frame.south west)
-- ([xshift=-1mm,yshift=-1mm]frame.north west)
[sharp corners]-- cycle;
}
,interior engine=empty,
},
fonttitle=\bfseries,
title=Remark \therem\; {#2},#1}



\newtcolorbox[use counter=df,number within=section]{df}[2][]{enhanced,
colbacktitle=\blue,
boxrule=0.2mm,
attach boxed title to top left={xshift=1cm,yshift*=1mm-\tcboxedtitleheight},
varwidth boxed title*=-3cm,
boxed title style={frame code={
\path[fill=tcbcolback!30!black]
([yshift=-1mm,xshift=-1mm]frame.north west)
arc[start angle=0,end angle=180,radius=1mm]
([yshift=-1mm,xshift=1mm]frame.north east)
arc[start angle=180,end angle=0,radius=1mm];
\path[left color=tcbcolback!60!black,right color=tcbcolback!60!black,middle color=tcbcolback!60!black]
([xshift=-2mm]frame.north west) -- ([xshift=2mm]frame.north east)
[rounded corners=1mm]-- ([xshift=1mm,yshift=-1mm]frame.north east)
-- (frame.south east) -- (frame.south west)
-- ([xshift=-1mm,yshift=-1mm]frame.north west)
[sharp corners]-- cycle;
}
,interior engine=empty,
},
fonttitle=\bfseries,
title=Definition \thedf\; {#2},#1}

\newtcolorbox[use counter=thm,number within=section]{thm}[2][]{
enhanced,
%skin=enhancedlast jigsaw,
attach boxed title to top left={xshift=-4mm,yshift=-0.5mm},
fonttitle=\bfseries\sffamily,varwidth boxed title=0.7\linewidth,
colbacktitle=lblue!60!black,colframe=lblue!60!black,
%interior style={top color=lblue1,bottom color=lred1},
boxed title style={empty,arc=0pt,outer arc=0pt,boxrule=0pt},
underlay boxed title={
\fill[lblue!80!black] (title.north west) -- (title.north east)
-- +(\tcboxedtitleheight-1mm,-\tcboxedtitleheight+1mm)
-- ([xshift=4mm,yshift=0.5mm]frame.north east) -- +(0mm,-1mm)
-- (title.south west) -- cycle;
\fill[lblue2] ([yshift=-0.5mm]frame.north west)
-- +(-0.4,0) -- +(0,-0.3) -- cycle;
\fill[lblue2] ([yshift=-0.5mm]frame.north east)
-- +(0,-0.3) -- +(0.4,0) -- cycle; 
},
title= Theorem \thethm\; {#2},#1}

%\thetcbcounter

\newtcolorbox[use counter=example,number within=section]{example}[1][]{%
empty,title={Example \theexample},attach boxed title to top left,
boxed title style={empty,size=minimal,toprule=2pt,top=4pt,
overlay={\draw[lblue!60!black,line width=2pt]
([yshift=-1pt]frame.north west)--([yshift=-1pt]frame.north east);}},
coltitle=lblue!60!black,fonttitle=\bfseries,
before=\par\medskip\noindent,parbox=false,boxsep=0pt,left=0pt,right=3mm,top=4pt,
breakable,pad at break*=0mm,vfill before first,
overlay unbroken={\draw[lblue,line width=1pt]
([yshift=-1pt]title.north east)--([xshift=-0.5pt,yshift=-1pt]title.north-|frame.east)
--([xshift=-0.5pt]frame.south east)--(frame.south west); },
overlay first={\draw[lblue,line width=1pt]
([yshift=-1pt]title.north east)--([xshift=-0.5pt,yshift=-1pt]title.north-|frame.east)
--([xshift=-0.5pt]frame.south east); },
overlay middle={\draw[lblue,line width=1pt] ([xshift=-0.5pt]frame.north east)
--([xshift=-0.5pt]frame.south east); },
overlay last={\draw[lblue,line width=1pt] ([xshift=-0.5pt]frame.north east)
--([xshift=-0.5pt]frame.south east)--(frame.south west);},#1%
}

\aliascntresetthe{df}
\aliascntresetthe{rem}
\aliascntresetthe{thm}
\aliascntresetthe{example}
\aliascntresetthe{exercise}

\newcommand*{\dfautorefname}{Definition}
\newcommand*{\remautorefname}{Remark}
\newcommand*{\thmautorefname}{Theorem}
\newcommand*{\exampleautorefname}{Example}
\newcommand*{\exerciseautorefname}{Exercise}
%\renewcommand*{\sectionautorefname}{Section}
%\renewcommand*{\subsectionautorefname}{Section}

\newtcolorbox{solution}{
empty,
breakable,left=5mm,right=5mm,
borderline west={1mm}{0pt}{\blue!60!black},
borderline east={1mm}{0pt}{\blue!60!black},
title= {Solution:\:\:},
attach title to upper,
fonttitle=\bfseries,
coltitle = black,
}

\newcommand\encircle[1]{%
  \tikz[baseline=(X.base)] 
    \node (X) [draw, shape=circle, inner sep=0] {\strut #1};}

\renewcommand*\contentsname{Contents}

\newcommand{\sitem}[1]{\begin{itemize}[noitemsep,nolistsep] #1 \end{itemize}}
\newcommand{\senum}[1]{\begin{enumerate}[noitemsep,nolistsep] #1 \end{enumerate}}
\newcommand{\tom}[1]{({\bf Tom:} \emph{#1})}
\newcommand{\defn}[1]{\begin{center} \emph{#1} \end{center}}
\newcommand{\pic}[2][0.7]{
\begin{center}
	\includegraphics[width=#1\textwidth]{images/#2}
\end{center}
}

%setting QR links
\newcommand{\qrlink}[2][]{\marginnote{\qrcode[height=0.9\marginparwidth,#1]{#2}}}%[-0.5cm]}



%\newcommand{\degs}{\degree}
\newcommand{\degs}{^\circ}

%\AtBeginDocument{%
% \abovedisplayshortskip=0pt
% \belowdisplayskip=0pt
% \belowdisplayshortskip=0pt
%}

\makeatletter
\newcommand\footnoteref[1]{\protected@xdef\@thefnmark{\ref{#1}}\@footnotemark}
\makeatother

\newcommand{\caution}{\begin{tikzpicture}[scale=1,baseline=0cm] %\scriptsize
	\path[rounded corners,draw=\yellow!50!\red,fill=none,thick]  (-0.25,0) -- (0.25,0) -- (0,0.433) -- cycle;
	\path[rounded corners=0.5mm,draw=none,fill=\yellow!60,thick]  (-0.19,0.04) -- (0.19,0.04) -- (0,0.37) -- cycle;
	\draw[draw=\yellow!20!\red,fill=none,thick]  (0,0.07) -- (0,0.1)  (0,0.12) -- (0,0.28);
\end{tikzpicture}}

\newcommand{\question}{
\tikz[baseline=(wi.base)]{
\node[fill=black,rotate=45,inner sep=.1ex,text height=1.8ex,text width=1.8ex] {};
\node[font=\color{white}] (wi) {?};
}	
}

\usepackage{xcolor}

\newcommand{\hl}[1]{%
  \colorbox{\red!50}{#1}}

%%
\usepackage{pgfplots,mathtools}
\pgfplotsset{compat=1.7}
\pgfmathdeclarefunction{gauss}{2}{\pgfmathparse{1/(#2*sqrt(2*pi))*exp(-((x-#1)^2)/(2*#2^2))}%
}

%%Importing tutorial questions
\usepackage{catchfilebetweentags}
\newtoks\temptoken

\usepackage[lastexercise,answerdelayed]{exercise}  
\renewcommand{\QuestionNB}{(\alph{Question})\ }
%\renewcommand{\subQuestionNB}{\roman{subQuestion})}
\renewcommand{\ExerciseHeader}{\textbf{\theExercise.} \hspace{-4mm} \textbf{\ExerciseHeaderTitle}}
\renewcommand{\AnswerHeader}{\textbf{\theExercise.}  \hspace{-4mm} \textbf{\ExerciseHeaderTitle}}
\setlength{\ExerciseSkipBefore}{-0.8ex}
\setlength{\AnswerSkipBefore}{-3ex}
\setlength{\QuestionBefore}{-1ex}

\renewcounter{Exercise}[subsection]% Reset counter.
\renewcommand{\theExercise}{\thesubsection.\arabic{Exercise}}%

\newcommand{\multicolq}[2][3]{
\begin{multicols*}{#1}
#2
\EndCurrentQuestion
\end{multicols*}
}

\newcommand{\multicolqq}[2][3]{
\begin{multicols*}{#1}
#2
\EndCurrentsubQuestion
\end{multicols*}
}

\newcommand{\tutorial}[1]{
\subsubsection*{Tutorial Exercises for \thesubsection} 

\medskip 

\ExecuteMetaData[\tutorialsheetlocation]{#1}

\vspace*{-1ex} \hspace*{\fill} \hyperref[A:answers]{(Answers)}
}

%%%%%%%Alex's Stuff%%%%%%%%%%%%

\newcommand{\bbA}{{\mathbb A}}
\newcommand{\bbB}{{\mathbb B}}
\newcommand{\bbC}{{\mathbb C}}
\newcommand{\bbD}{{\mathbb D}}
\newcommand{\bbE}{{\mathbb E}}
\newcommand{\bbF}{{\mathbb F}}
\newcommand{\bbG}{{\mathbb G}}
\newcommand{\bbH}{{\mathbb H}}
\newcommand{\bbI}{{\mathbb I}}
\newcommand{\bbJ}{{\mathbb J}}
\newcommand{\bbK}{{\mathbb K}}
\newcommand{\bbL}{{\mathbb L}}
\newcommand{\bbM}{{\mathbb M}}
\newcommand{\bbN}{{\mathbb N}}
\newcommand{\bbO}{{\mathbb O}}
\newcommand{\bbP}{{\mathbb P}}
\newcommand{\bbQ}{{\mathbb Q}}
\newcommand{\bbR}{{\mathbb R}}
\newcommand{\bbS}{{\mathbb S}}
\newcommand{\bbT}{{\mathbb T}}
\newcommand{\bbU}{{\mathbb U}}
\newcommand{\bbV}{{\mathbb V}}
\newcommand{\bbW}{{\mathbb W}}
\newcommand{\bbX}{{\mathbb X}}
\newcommand{\bbY}{{\mathbb Y}}
\newcommand{\bbZ}{{\mathbb Z}}

\newcommand{\cA}{{\mathcal A}}
\newcommand{\cB}{{\mathcal B}}
\newcommand{\cC}{{\mathcal C}}
\newcommand{\cD}{{\mathcal D}}
\newcommand{\cE}{{\mathcal E}}
\newcommand{\cF}{{\mathcal F}}
\newcommand{\cG}{{\mathcal G}}
\newcommand{\cH}{{\mathcal H}}
\newcommand{\cI}{{\mathcal I}}
\newcommand{\cJ}{{\mathcal J}}
\newcommand{\cK}{{\mathcal K}}
\newcommand{\cL}{{\mathcal L}}
\newcommand{\cM}{{\mathcal M}}
\newcommand{\cN}{{\mathcal N}}
\newcommand{\cO}{{\mathcal O}}
\newcommand{\cP}{{\mathcal P}}
\newcommand{\cQ}{{\mathcal Q}}
\newcommand{\cR}{{\mathcal R}}
\newcommand{\cS}{{\mathcal S}}
\newcommand{\cT}{{\mathcal T}}
\newcommand{\cU}{{\mathcal U}}
\newcommand{\cV}{{\mathcal V}}
\newcommand{\cW}{{\mathcal W}}
\newcommand{\cX}{{\mathcal X}}
\newcommand{\cY}{{\mathcal Y}}
\newcommand{\cZ}{{\mathcal Z}}

\newcommand{\fra}{{\mathfrak a}}
\newcommand{\frb}{{\mathfrak b}}
\newcommand{\frc}{{\mathfrak c}}
\newcommand{\frd}{{\mathfrak d}}
\newcommand{\fre}{{\mathfrak e}}
\newcommand{\frf}{{\mathfrak f}}
\newcommand{\frg}{{\mathfrak g}}
\newcommand{\frh}{{\mathfrak h}}
\newcommand{\fri}{{\mathfrak i}}
\newcommand{\frj}{{\mathfrak j}}
\newcommand{\frk}{{\mathfrak k}}
\newcommand{\frl}{{\mathfrak l}}
%\renewcommand{\frm}{{\mathfrak m}}
\newcommand{\frn}{{\mathfrak n}}
\newcommand{\fro}{{\mathfrak o}}
\newcommand{\frp}{{\mathfrak p}}

\newcommand{\frr}{{\mathfrak r}}
\newcommand{\frs}{{\mathfrak s}}
\newcommand{\frt}{{\mathfrak t}}
\newcommand{\fru}{{\mathfrak u}}
\newcommand{\frv}{{\mathfrak v}}
\newcommand{\frw}{{\mathfrak w}}
\newcommand{\frx}{{\mathfrak x}}
\newcommand{\fry}{{\mathfrak y}}
\newcommand{\frz}{{\mathfrak z}}


\newcommand{\frA}{{\mathfrak A}}
\newcommand{\frB}{{\mathfrak B}}
\newcommand{\frC}{{\mathfrak C}}
\newcommand{\frD}{{\mathfrak D}}
\newcommand{\frE}{{\mathfrak E}}
\newcommand{\frF}{{\mathfrak F}}
\newcommand{\frG}{{\mathfrak G}}
\newcommand{\frH}{{\mathfrak H}}
\newcommand{\frI}{{\mathfrak I}}
\newcommand{\frJ}{{\mathfrak J}}
\newcommand{\frK}{{\mathfrak K}}
\newcommand{\frL}{{\mathfrak L}}
\newcommand{\frM}{{\mathfrak M}}
\newcommand{\frN}{{\mathfrak N}}
\newcommand{\frO}{{\mathfrak O}}
\newcommand{\frP}{{\mathfrak P}}
\newcommand{\frQ}{{\mathfrak Q}}
\newcommand{\frR}{{\mathfrak R}}
\newcommand{\frS}{{\mathfrak S}}
\newcommand{\frT}{{\mathfrak T}}
\newcommand{\frU}{{\mathfrak U}}
\newcommand{\frV}{{\mathfrak V}}
\newcommand{\frW}{{\mathfrak W}}
\newcommand{\frX}{{\mathfrak X}}
\newcommand{\frY}{{\mathfrak Y}}
\newcommand{\frZ}{{\mathfrak Z}}

\newcommand{\lra}{\longrightarrow}
\newcommand{\ra}{\rightarrow}
\newcommand{\Ra}{\Rightarrow}
\newcommand{\La}{\Leftarrow}
\newcommand{\Lra}{\Leftrightarrow}
\newcommand{\hra}{\hookrightarrow}
\newcommand{\eps}{\varepsilon}
\newcommand{\Vxi}{V_{\cU, \varepsilon}(\xi)}
\newcommand{\Veta}{V_U(\eta)}



\newcommand{\xh}{\hat{x}}
\newcommand{\mata}[1]{\left(\begin{array}{cc}#1\end{array}\right)}
\newcommand{\asu}{\mathfrak{su}}
\newcommand{\sSU}{\mathsf{SU}}
\newcommand{\sPSU}{\mathsf{P SU}}
\newcommand{\CC}{\mathcal{C}}
\newcommand{\CCC}{\mathscr{C}}
\newcommand{\CCL}{\mathscr{L}}
%\newcommand{\CD}{\mathcal{D}}
\newcommand{\CCD}{\mathscr{D}}
\newcommand{\CCDb}{\bar{\mathscr{D}}}
\newcommand{\CF}{\mathcal{F}}
\newcommand{\CCF}{\mathscr{F}}
\newcommand{\CG}{\mathcal{G}}
\newcommand{\CCG}{\mathscr{G}}
\newcommand{\CH}{\mathcal{H}}
\newcommand{\CCH}{\mathscr{H}}
\newcommand{\CI}{\mathcal{I}}
\newcommand{\CJ}{\mathcal{J}}
\newcommand{\CK}{\mathcal{K}}
\newcommand{\sU}{\mathsf{U}}
\newcommand{\sSL}{\mathsf{SL}}
\newcommand{\sGL}{\mathsf{GL}}
\newcommand{\sMat}{\mathsf{Mat}}
\newcommand{\sDiag}{\mathsf{Diag}}
\newcommand{\sO}{\mathsf{O}}
\newcommand{\de}{\mathrm{e}}     			% Euler's number
\newcommand{\rk}{{\rm rk}}     			% imaginary unit
\newcommand{\di}{\mathrm{i}}     			% imaginary unit
\newcommand{\sSO}{\mathsf{SO}}
\newcommand{\sS}{\mathsf{S}}
\newcommand{\sSpin}{\mathsf{Spin}}
\newcommand{\sPin}{\mathsf{Pin}}
\newcommand{\Det}{\hbox{Det}}
\newcommand{\hbo}{\hbox to 1 true cm {\hfill } }
\newcommand{\RZ}{\mathds{Z}}     			% ring of integers
%\renewcommand{\span}{{\rm span}}
\newcommand{\tr}{\hbox{tr}}
\newcommand{\Tr}{\hbox{Tr}}
\newcommand{\slc}{$\sSL(2,\mathbb{C})$}
\newcommand{\dalpha}{{\dot\alpha}}
\newcommand{\dbeta}{{\dot\beta}}
%\newcommand{\eps}{{\varepsilon}}			% antisymmetric tensors
\renewcommand{\sup}[2]{\{\![#1,#2]\!\}}
\newcommand{\vectt}[2]{\left(\begin{array}{c} #1 \\ #2 \end{array}\right)}
\newcommand{\vecttt}[3]{\left(\begin{array}{c} #1 \\ #2 \\ #3\end{array}\right)}
\newcommand{\vectttt}[4]{\left(\begin{array}{c} #1 \\ #2 \\ #3 \\ #4\end{array}\right)}
\newcommand{\vecttdt}[3]{\left(\begin{array}{c} #1 \\ #2 \\ \vdots \\ #3\end{array}\right)}
\newcommand{\matrixbig}[3]{\left(\begin{array}{cccc} #1_{11} & #1_{12} & \ldots & #1_{1#3} \\
		#1_{21} & #1_{22} & \ldots & #1_{2#3} \\
		\vdots & \vdots &  & \vdots \\
		#1_{#2 1} & #1_{#2 2} & \ldots & #1_{#2 #3}\end{array}\right)}
\newcommand{\matrixsmall}[3]{\left(\begin{array}{ccc} #1_{11} & \ldots & #1_{1#3} \\
		\vdots &   & \vdots \\
		#1_{#2 1} & \ldots & #1_{#2 #3}\end{array}\right)}
\newcommand{\vectbig}[2]{\left(\begin{array}{c} #1_1 \\ #1_2 \\ \vdots \\ #1_#2\end{array}\right)}
\newcommand{\ggama}{{\boldsymbol\gamma}}
\newcommand{\un}[1]{\underline{#1}}
\newcommand{\Eins}{\mathdsss{1}}
%\newcommand{\dd}{\mathrm{d}}     			% total differential
\newcommand{\bracket}[2]{\bra{#1}\,#2\rangle} % Dirac inner product
\newcommand{\bra}[1]{\langle\,#1\,|}          % Dirac bra
\newcommand{\ket}[1]{|\,#1\,\rangle}          % Dirac Ket
\newcommand{\linv}{{\triangle}^{-1}}      % inv. Laplacian
\newcommand{\ud}{\mathrm{d}}
%\newcommand{\remark}[1]{}     				% remark
\newcommand{\ewith}{~~~\mbox{with}~~~}
\newcommand{\eand}{~~~\mbox{and}~~~}
\newcommand{\eor}{~~~\mbox{or}~~~}
\newcommand{\nextyear}[1]{}     				% remark
\newcommand{\pathD}{\mathcal{D}}
\newcommand{\e}{\mathrm{e}}		% exponential e
\newcommand{\x}{\boldsymbol{x}}		% vector x
\newcommand{\NN}{\mathds{N}}     			% set of natural numbers
\newcommand{\unit}{\mathds{1}}   			% identity map/matrix
\newcommand{\y}{\boldsymbol{y}}		% vector y
\newcommand{\z}{\boldsymbol{z}}		% vector z
\newcommand{\av}{\vec{a}}		% vector v
\newcommand{\bv}{\vec{b}}		% vector v
\newcommand{\cv}{\vec{c}}		% vector v
\newcommand{\ev}{\vec{e}}		% vector v
\newcommand{\fv}{\vec{f}}	
\newcommand{\nv}{\vec{0}}
\newcommand{\pv}{\vec{p}}		% vector v
\newcommand{\rv}{\vec{r}}		% vector v
\newcommand{\sv}{\vec{s}}		% vector v
\newcommand{\uv}{\vec{u}}		% vector v
\newcommand{\vv}{\vec{v}}		% vector v
\newcommand{\wv}{\vec{w}}		% vector v
\newcommand{\xv}{\vec{x}}		% vector v
\newcommand{\yv}{\vec{y}}		% vector v
\newcommand{\zv}{\vec{z}}	
\renewcommand{\u}{\boldsymbol{u}}
\newcommand{\p}{\boldsymbol{p}}		% vector p
\newcommand{\q}{\boldsymbol{q}}		% vector q
\renewcommand{\r}{\boldsymbol{r}}		% vector q
%\newcommand{\E}{\boldsymbol{E}}
\newcommand{\A}{\boldsymbol{A}}
\newcommand{\J}{\boldsymbol{J}}
%\newcommand{\Q}{\mathcal{Q}}
\newcommand{\FR}{\mathds{R}} 
\newcommand{\FZ}{\mathds{Z}} 
\newcommand{\FN}{\mathds{N}}     			% field of real numbers
\newcommand{\FC}{\mathds{C}}     			% field of complex numbers
\newcommand{\CN}{\mathcal{N}}   
\newcommand{\CP}{\mathcal{P}}   	
\newcommand{\CB}{\mathcal{B}} 
\newcommand{\CL}{\mathcal{L}}   		% field of real numbers
%\newcommand{\eolec}[2]{}
%\newcommand{\eolec}[2]{\marginpar{~~~~~#1\\[-0.3cm]\rule{50pt}{0.4pt}\\\phantom{~~~~~}#2}}
%\newcommand{\tut}[1]{\marginpar{Tutorial #1\\[-0.3cm]\rule{50pt}{0.4pt}\\[-0.45cm]\rule{50pt}{0.4pt}}}

\usepackage{dsfont}

%\newcommand{\hey}[1]{{\color{TealBlue}(#1)}}

\newcommand{\elt}[1]{\stackrel{\footnotesize{#1}}{\rightsquigarrow}}
\newcommand{\melt}[1]{\stackrel{\begin{smallmatrix}#1\end{smallmatrix}}{\rightsquigarrow}}
